\scanpage{77}\mspage{\unsure{64}}
% archive sheet number: 76
% Paragraph continues from the preceding sheet.
ce foncteur est adj. à g. de $w_W^*$ !

Comme on a évidemment
\[
u_!^W\simeq v_!^W w_!^W,\qquad u_W^*\simeq \deleted{\unsure{W}}w_W^*\circ v_W^{\unsure{'|*}}
\]
on en conclut que $u_!^W$ est bien adjoint
à gauche de $u_W^*$, \unsure{q ed.|q.e.d.}

\uline{Corollaire} 3.9. Soit $(M,W,C)$ une catégorie
à cofibrations. Alors le prédérivateur
$\mathbb{D}_W$, \deleted{\unsure{satisf}} restreint à la catégorie \deleted{\unsure{des}}
$\mathrm{Diag}_0$ formée des cat. \inserted{$I$} \uline{ordonnées} \deleted{\unsure{$I$}}
artiniennes et telles que $\mathrm{Esp}\,I$
soit noethérien, satisfait aux
axiomes Der 1 (\uline{décomposition}), Der 2
(\uline{\unsure{localisation|localisabilité}}), Der 3 (\uline{objet final}),
Der 4 bis (\uline{existence des $u_!$}), Der 5 bis
(mode de \uline{calcul des $u_!$}), Der 6
(\uline{exactitude des} carrés \inserted{$W$-}\uline{cocartésiens}).

Si $C$, $TC=C\cap W$ sont stables par composition
transfinie, on a le $\unsure{=}$ résultat
sur la catégorie $\mathrm{Diag}$ formée des
cat $I$ ordonnées artiniennes sans plus.
